% Isolated proposed B167, checkpoint35. Earlier source bytes remain unchanged.
\subsection*{An entire compact contact family in the remaining kernel}
\statement{B.167}{Proved}{The actual \(m\)-orbit of \(D_s\) is a free profinite continuous family, embeds in the whole contact presentation without relations, and has a strictly different topology from its discrete infinite-rank kernel}
\label{stmt:B167}
Retain the actual marked pair
\[
 H_D=\langle m,D_s\rangle\le\Gamma_0
\]
of B.132. It is free of rank two and has its full induced profinite
topology in the actual group with the entire node filling.
Write \(\mathcal H_D=\overline{H_D}\cong\widehat F(m,D_s)\).
Use the filled one-family projection, the actual source section, and
the closed groups of B.150, B.154, and B.164:
\[
 \begin{gathered}
 \widehat q_a:\widehat{\Gamma_0}\longrightarrow\widehat{\Gamma_a},
 \qquad
 \widehat J=\widehat{\langle A,m,v\mid[A,m]=1\rangle},
 \qquad
 \widehat\sigma:\widehat J\longrightarrow\widehat{\Gamma_0},\\
 \mathcal K=\ker\widehat q_a,\qquad
 \mathcal I=\overline{H_{\mathrm{src}}}\cap\mathcal K,\qquad
 \mathcal U_D=\mathcal H_D\cap\mathcal K.
 \end{gathered}
\]
Then \(\mathcal U_D\) is the free profinite group on the compact
generating space \(\widehat{\mathbb Z}\), with the exact continuous
marked family
\begin{equation}\label{eq:B167-family}
 d_z=m^zD_sm^{-z}\qquad(z\in\widehat{\mathbb Z}).
\end{equation}
It is contained in \(\mathcal I\le\mathcal K^\delta\).
For the entire continuous contact presentation of B.164, put
\[
 W=F_{\mathrm{prof}}\bigl(\widehat J\times\{B,D,E\}\bigr),
 \qquad \operatorname{ev}:W\twoheadrightarrow\mathcal I,
 \qquad\mathcal R_{\mathrm{src}}=\ker\operatorname{ev}.
\]
The closed subgroup \(W_D\) generated by \(e_{m^z,D}\) is a retract
of \(W\), and evaluation identifies it with \(\mathcal U_D\).
In particular
\begin{equation}\label{eq:B167-restriction}
                 \mathcal R_{\mathrm{src}}\cap W_D=1.
\end{equation}
The actual compact family is not the full completion of the discrete
kernel \(U_{\mathrm{disc}}=\ker(H_D\to\mathbb Z_m)\):
the canonical map
\begin{equation}\label{eq:B167-not-completion}
        \widehat{U_{\mathrm{disc}}}\twoheadrightarrow\mathcal U_D
\end{equation}
is surjective and noninjective.

\begin{proof}
\emph{The complete continuous-family construction.}
For each \(n\ge1\), let
\[
 F_n=\widehat{F(d_i:i\in\mathbb Z/n)}.
\]
If \(n\mid n'\), send \(d_i\) at level \(n'\) to \(d_{i\bmod n}\).
These are surjective continuous homomorphisms. Put
\(U=\varprojlim_nF_n\). The compatible generators give a continuous
map \(z\mapsto d_z\) from \(\widehat{\mathbb Z}\) to \(U\).
Every continuous map from \(\widehat{\mathbb Z}\) to a finite group
is constant on the cosets of \(n\widehat{\mathbb Z}\) for some \(n\):
its finitely many clopen fibers admit one common such refinement.
The finite-rank universal property then extends the map through
\(F_n\), hence through \(U\).
The family topologically generates \(U\), since it generates each
coordinate and every finite set of coordinates has a common refinement.
This proves uniqueness. Compatible finite-target extensions give the
same universal property for every profinite target. Thus \(U\) is
exactly the free profinite group on this compact space; arbitrary
discontinuous assignments on all integer labels are not included.

Translation acts on \(F_n\) through \(\mathbb Z/n\), compatibly in \(n\).
It gives a jointly continuous action of \(\widehat{\mathbb Z}\) on
\(U\): each coordinate action is continuous on
\((\mathbb Z/n)\times F_n\). In \(U\rtimes\widehat{\mathbb Z}\), call
the procyclic generator \(m\).
The free profinite map from \(\widehat F(m,D)\) sends \(m\) to this
element and \(D\) to \(d_0\). Conversely the continuous family
\[
                 z\longmapsto m^zDm^{-z}
                 \quad\text{in }\widehat F(m,D)
\]
gives a homomorphism \(U\to\widehat F(m,D)\). Its translation
equivariance and the procyclic section give a homomorphism from
\(U\rtimes\widehat{\mathbb Z}\). Both composites fix \(m,d_0\);
conjugation by every \(m^z\) shows that they fix every \(d_z\).
They are inverse continuous isomorphisms. In particular
\begin{equation}\label{eq:B167-split}
 \widehat F(m,D)
   \cong U\rtimes\widehat{\mathbb Z},
 \qquad
 U=\ker\bigl(\widehat F(m,D)\longrightarrow\widehat{\mathbb Z}\bigr).
\end{equation}

\emph{The actual marked subgroup.}
B.132 supplies the full-topology embedding of the whole pair in
\(\widehat{\Gamma_0}\), with all forty ordinary node seeds, both
sheets, and every signed source translate retained.
The filled map sends \(D_s\) to \(1\) and \(m\) to the indicated
generator of \(\widehat J\). The \(m\)-exponent character is a
continuous left inverse on this procyclic subgroup, so it is
\(\widehat{\mathbb Z}\). Therefore the restriction of
\(\widehat q_a\) to \(\mathcal H_D\) is precisely the projection in
\eqref{eq:B167-split}. This proves \eqref{eq:B167-family} and the
claimed actual intersection. The two generators are fixed source
elements. Their closure, and hence \(\mathcal U_D\), is pointwise
\(\delta\)-fixed. Its elements lie in the source closure and in
\(\mathcal K\); B.154 therefore places them in the actual
\(\mathcal I\), without replacing it by an ambient normal closure.

\emph{Restriction of the entire relation kernel.}
The character
\(\chi_m:\widehat J\to\widehat{\mathbb Z}\) has
\(\chi_m(m)=1\) and \(\chi_m(A)=\chi_m(v)=0\).
The continuous assignments
\[
 e_{j,D}\longmapsto d_{\chi_m(j)},\qquad
 e_{j,B},e_{j,E}\longmapsto1
\]
give a homomorphism \(r:W\to U\).
Conversely \(d_z\mapsto e_{m^z,D}\) gives \(i:U\to W\).
The composite \(ri\) fixes every \(d_z\), hence is the identity.
Thus \(i\) is a compact embedding with retract image \(W_D\).
Evaluation sends \(i(d_z)\) to the literal element
\(m^zD_sm^{-z}\). The preceding actual pair comparison says that
\(\operatorname{ev}i\) is injective, proving
\eqref{eq:B167-restriction}.
The retraction here is a retraction of \(W\); it is not asserted
to descend to \(\mathcal I\), and no free-product position in
\(\mathcal I\) is inferred.

\emph{Exact cohomological placement.}
B.164 gives an injective continuous connecting map
\[
 \partial:\mathcal K^\delta\longrightarrow
 H^1(\langle\delta\rangle,\mathcal N)
\]
whose image is the basepoint fiber of the map to
\(H^1(\langle\delta\rangle,\mathcal K_a)\).
Its restriction to \(\mathcal U_D\) is a homeomorphism onto a
closed subspace of \(\partial(\mathcal I)\).
The group law may be transported from \(\mathcal U_D\) to this
subspace; no group law on the entire nonabelian cohomology space
is thereby asserted. In particular this does not identify
\(\partial(\mathcal I)\) with the entire basepoint fiber.

\emph{The negative discrete-topology statement.}
Ordinary Schreier rewriting gives the basis
\[
 d_j=m^jD_sm^{-j}\quad(j\in\mathbb Z)
\]
of \(U_{\mathrm{disc}}\). Define a \(C_2\)-character \(\eta\) by
\(\eta(d_0)=1\) and \(\eta(d_j)=0\) for \(j\ne0\).
For any homomorphism from the actual ambient \(\Gamma_0\) to a
finite group, let \(n\ge1\) be the order of the image of \(m\).
Then \(d_nd_0^{-1}\) lies in its kernel, whereas
\(\eta(d_nd_0^{-1})=1\).
No such ambient finite kernel refines \(\ker\eta\).
Thus \(U_{\mathrm{disc}}\) does not have its full abstract
profinite topology in \(\Gamma_0\), despite the full topology
of its containing two-generator group \(H_D\).

Integer labels are dense in the compact family
\eqref{eq:B167-family}, so the canonical map
\eqref{eq:B167-not-completion} is onto.
If it were injective, it would be a compact homeomorphism and
\(\eta\) would become continuous on \(\mathcal U_D\).
Every continuous finite quotient kernel of this closed subgroup
contains its intersection with some open normal subgroup of
\(\widehat{\Gamma_0}\): take an ambient open neighborhood inside
the subgroup kernel and then an open normal neighborhood.
Restriction to \(\Gamma_0\) would give the forbidden refinement.
Hence \eqref{eq:B167-not-completion} is noninjective.
This is a kernel-containment argument; extension to the same
finite codomain is not required.
\end{proof}

\begingroup\footnotesize\raggedright
The full prior proof and exact input bindings are
\path{work_CE/compact_single_contact_relation_restriction34.txt}
and its JSON companion. The complete independent review is
\path{work_DE/review_compact_single_contact_relation_restriction34.md}.
All topology statements above refer to the actual compact subgroup
and continuous finite targets. They do not say that every abstract
finite-index subgroup of an infinitely generated profinite group is
open. The full three-contact relation kernel, the contact-generation
equality, and the separate original B.75 gate remain undetermined.
\par\endgroup
